%%%% ABSTRACT
%%
%% Versão do resumo para idioma de divulgação internacional.

\begin{abstractutfpr}%% Ambiente abstractutfpr
The detection of insect pests in agricultural crops is fundamental for Integrated Pest Management (IPM); however, strictly spatial vision-based approaches frequently fail in critical camouflage scenarios and with small-scale targets. This work proposes and evaluates a hybrid Deep Learning architecture based on the U-Net network, extended to process spatiotemporal data. To break the visual ambiguity imposed by mimicry (where the color and texture of the insect statistically blend with the leaf background), the model incorporates kinematic estimation through Farneback's Dense Optical Flow, resulting in 6-channel input tensors (RGB and flow components). To enable training without exhaustive manual annotation, a weak supervision pipeline was implemented combining CLAHE adaptive equalization and Gaussian Mixture Models (MOG2). The experimental results proved that the traditional spatial U-Net network fails in the extreme green-on-green camouflage scenario due to class imbalance and lack of contrast. The application of data augmentation partially mitigated overfitting, but only the early fusion of motion signatures was able to consistently segment the anatomy of the specimens in all scenarios. It is concluded that incorporating temporal bias is highly effective in regularizing convolutional networks for small-scale biological detection problems, offering a robust alternative for autonomous leaf diagnostic systems.
\end{abstractutfpr}
